\section{An explicit saving for a fixed two-coordinate transform}
\label{net:section}
\label{sec:network}

Put
\begin{equation}\label{net:C}
 C=\begin{pmatrix}a&b\\b&a\end{pmatrix},
 \qquad a=\frac{1+i}{2},\qquad b=\frac{1-i}{2}.
\end{equation}
The ordinary algorithm applies the $k$ factors of $C^{\otimes k}$
separately, using $O(k2^k)$ operations.  We obtain a smaller exponent
of $k$ by a fixed network from the complex construction of
\cite[Proposition~8 and Section~3.5]{motif}.  The network exchanges two banks and
restores arbitrary auxiliary values.  Changing the coordinate frame
along its wires then transforms \emph{every} bank and auxiliary array.
The total dimension of the frame changes is smaller than the cost of
transforming the arrays separately.  We give the complete finite
construction because this strict saving and the treatment of arbitrary
auxiliaries are the main inputs to the recursion.

\subsection{A scalar network with arbitrary auxiliary values}

Fix
\begin{equation}\label{net:parameters}
 h=100,\quad
 \mathcal T=\{T\subseteq\{1,\ldots,h\}:|T|=3\},\quad
 v=\binom h3=161700,\quad N_0=v^3,\quad m=h^3=10^6.
\end{equation}
There are two banks of scalar wires $X_d,Y_d$, indexed by
$d=(d_1,d_2,d_3)\in\mathcal T^3$.  Two triples are \emph{neighbors}
when they have even intersection; neighboring triples are distinct.
The network has three stages.  At stage $j$, fix the two indices other
than $d_j$ and operate on the two lists obtained by varying $d_j$.
There are $I_0=3v^2$ such invocations.  Each has its own auxiliary
wires, distinct from all other invocations: one side wire $A_{ST}$
for each \emph{ordered} neighboring pair $(S,T)$, and $h+1$ central
wires $c_1,\ldots,c_h,c_*$.

For one invocation, call the logical source and target banks
$x=(x_T)$ and $y=(y_S)$.  Define the fixed linear maps
\begin{align*}
 (Vx)_{ST}&=x_T,&
 (Gx)_j&=\sum_{T\ni j}x_T,& (Gx)_*&=\sum_Tx_T,\\
 (JA)_S&=-\sum_{T:\,T\text{ neighbors }S}
                  \frac{|S\cap T|-1}{2}A_{ST},&
 (Rc)_S&=\frac{\sum_{j\in S}c_j-c_*}{2}.
\end{align*}
The forward invocation consists of the following eight rows, in order:
\begin{equation}\label{net:schedule}
\begin{array}{c|l|l}
 \text{row}&\text{update}&\text{gate grouping}\\\hline
 0&y\gets y-JA&\text{one gate per target and its side wires}\\
 1&y\gets y-Rc&\text{one gate on the targets and center}\\
 2&A\gets A+Vx&\text{one gate per source and its side wires}\\
 3&c\gets c+Gx&\text{one gate on the sources and center}\\
 4&y\gets y+Rc&\text{one gate on the targets and center}\\
 5&y\gets y+JA&\text{one gate per target and its side wires}\\
 6&c\gets c-Gx&\text{one gate on the sources and center}\\
 7&A\gets A-Vx&\text{one gate per source and its side wires}.
\end{array}
\end{equation}
A gate is simply the indicated finite linear map on its wires.  Its
scalar implementation, including all additions and multiplications,
has a fixed finite cost.

\begin{lemma}\label{net:scalar}
For arbitrary initial auxiliary values, the forward invocation sends
$(x,y,A,c)$ to $(x,y+x,A,c)$.  Forward invocations from $X$ to $Y$
at stages $1$ and $3$, and the inverse forward invocation from $Y$
to $X$ at stage $2$, give
\[
 (X_d,Y_d)\longmapsto(-Y_d,X_d)
 \quad(d\in\mathcal T^3)
\]
and restore every auxiliary value.
\end{lemma}

\begin{proof}
The first two rows subtract $JA+Rc$ from $y$.  The middle rows add
$R(c+Gx)+J(A+Vx)$, and the last two restore $c$ and $A$.  The net
addition is $(RG+JV)x$.  The coefficient from $x_T$ to $y_S$ in
$RG$ is $(|S\cap T|-1)/2$.  It is $1$ for $S=T$, $0$ when the
intersection has size $1$, and is canceled by $JV$ when the
intersection has size $0$ or $2$.  Hence $RG+JV=\Id$.
The three stages therefore send
$(x,y)$ to $(x,y+x)$, then $(-y,y+x)$, then $(-y,x)$.
The auxiliary values cancel in each invocation separately.
\end{proof}

A triple has
\[
 d_{\mathrm{nbr}}=\binom{97}{3}+3\cdot97=147731
\]
neighbors: the two terms count intersections of size $0$ and $2$.
Consequently the number of physical wires is
\begin{equation}\label{net:wire-count}
 W=2N_0+I_0(vd_{\mathrm{nbr}}+101)
   =1873807244643542670000.
\end{equation}
The factor $vd_{\mathrm{nbr}}$ counts ordered pairs, as required by
the schedule.

\subsection{Frames turn binary dimensions into directional operations}

We now explain how to turn the scalar exchange into a tensor transform.
An edge of the scalar network connects consecutive touches of a physical
wire, including its source and sink.  We will assign a binary subspace
to each gate and terminal, with one common subspace for all wires touched
by a gate.  This subspace specifies a change of frame on the complex
array replacing each wire.  The construction below has two useful
properties: common frames commute with scalar gates, and moving between
suitable nested subspaces costs one directional operation per added or
removed binary dimension.

For this purpose use the binary space $\mathcal D=\Ftwo^m$ with its
ordinary dot product.  A subspace is \emph{nondegenerate} if its restricted
bilinear form has zero radical.  Its orthogonal complement and all
projections below are taken over $\Ftwo$; these are separate from the
complex spaces carrying the data.

For $z\in\Ftwo^m$, write $R_z$ for translation of array addresses:
$(R_zf)(x)=f(x+z)$.  The normalized Walsh matrix
\[
 J_m=2^{-m/2}\bigl((-1)^{x\cdot y}\bigr)_{x,y\in\Ftwo^m}
\]
satisfies $J_m^2=\Id$.  Conjugating the diagonal character
$(-1)^{z\cdot x}$ by $J_m$ gives $R_z$, by cancellation of
nontrivial binary characters.  For a nondegenerate label $U$, let
$P_U$ be projection onto $U$ along $U^\perp$, and set
\begin{equation}\label{net:frame}
 q_U(x)=\wt(P_Ux)\pmod4,\qquad
 \Phi_U=J_m\diag_x(i^{q_U(x)})J_m.
\end{equation}
Here $\wt$ is integer Hamming weight.  These matrices prove
identities; the algorithm does not execute Walsh transforms.
In particular $\Phi_0=\Id$ and $\Phi_{\mathcal D}=C^{\otimes m}$.

\begin{lemma}\label{net:edge}
Let $U,V\subseteq\mathcal D$ be nested nondegenerate subspaces.
Call the orthogonal complement of the smaller subspace inside the
larger their \emph{residual}.  Suppose it has an orthonormal binary
basis of size $\rho$.
Then $\Phi_V\Phi_U^{-1}$ is a product of exactly $\rho$
\emph{directional kernels} of the form $C_z=a\Id+bR_z$ or $C_z^{-1}$,
where $z\ne0$.  Each inverse requires one forward kernel and a
translation.
\end{lemma}

\begin{proof}
For binary vectors $x,y$,
\[
 \wt(x+y)=\wt(x)+\wt(y)-2|\supp(x)\cap\supp(y)|.
\]
Thus weight is additive modulo four for orthogonal vectors.
If $V=U\mathbin\perp E$ and $z_1,\ldots,z_\rho$ is an
orthonormal basis of $E$, then
\[
 q_V(x)-q_U(x)
   =\sum_{j=1}^{\rho}\wt(z_j)[z_j\cdot x]\pmod4,
\]
where brackets mean the integer representative $0$ or $1$.
Every $\wt(z_j)$ is odd, so its value modulo four is $1$ or $-1$.
For a decreasing edge the entire expression is negated.
The identity
$i^{[z\cdot x]}=a+b(-1)^{z\cdot x}$ and Walsh conjugation
give the stated kernels.  Finally $R_z^2=\Id$ and
$C_z^2=R_z$, whence $C_z^{-1}=R_zC_z$.
\end{proof}

To apply this calculation to the network, replace each scalar wire by
an array indexed by $\mathcal D$.  Give every gate a nondegenerate label
$U$, common to all wires it touches, and give each source and sink its
own label.  Insert $\Phi_V\Phi_U^{-1}$ along an edge from label $U$ to
label $V$, and apply the scalar gates pointwise.  A scalar gate mixes
its touched arrays with scalar coefficients, whereas their common
frame acts on the address coordinate.  These two operations commute.
The intermediate frames therefore cancel, leaving the scalar network
between its source and sink frames.

More generally, for an integer $f\ge1$, index every array by $f$ columns
of $m$ bits and use $\Phi_U^{\otimes f}$ as its frame.  The same
cancellation gives the complete linear map
\begin{equation}\label{net:telescoping}
 \diag_{\mathrm{sinks}}(\Phi_U^{\otimes f})\,
 \mathcal S\,
 \diag_{\mathrm{sources}}((\Phi_U^{-1})^{\otimes f}),
\end{equation}
where $\mathcal S$ is the pointwise signed bank exchange of
Lemma~\ref{net:scalar}, acting as the identity on auxiliary roles.
This is an identity on arbitrary values in every physical role.
The remaining task is to choose labels whose terminal frame ratios
give the desired transform and whose total residual dimension is less
than $Wm$.  Lemma~\ref{net:edge} will then convert that dimension
saving into fewer directional operations.

\subsection{Binary labels and the dimension budget}

Identify $\mathcal D$ with $D^{\otimes3}$, where $D=\Ftwo^h$;
the tensor dot product is the ordinary dot product on $\Ftwo^m$.
For $T\in\mathcal T$, let $t_T$ be its indicator.
Then $t_T\cdot t_T=1$, and neighboring indicators are orthogonal.
Put
\[
 u_d=t_{d_1}\otimes t_{d_2}\otimes t_{d_3},\qquad
 U_d=\langle u_d\rangle.
\]
Each $u_d$ has norm one, so $U_d$ is nondegenerate.

Start with the terminal labels.  Sources have label $0$, except that
$X_d$ has label $U_d$; sinks have label $\mathcal D$, except that
$Y_d$ has label $U_d^\perp$.  A route with terminal labels $0$ and
$\mathcal D$ has frame ratio $C^{\otimes m}$.  The exceptional
route $X_d\to Y_d$ will have the same ratio up to a translation, as
we verify below.  The total increase of terminal dimensions is
$Wm-2N_0$.  Thus the margin below $Wm$ is $2N_0$ before we
charge twice the dimensions lost at any internal decreases.

At an invocation in stage $j$, write
\[
 E=D^{\otimes(j-1)},\qquad
 P=\Big\langle\bigotimes_{l<j}t_{d_l}\Big\rangle,
 \qquad B=P^\perp\subset E,\qquad
 Q=\Big\langle\bigotimes_{l>j}t_{d_l}\Big\rangle.
\]
An empty tensor product is the full one-dimensional binary space.
Thus $E=B\mathbin\perp P$, with $B=0$ in stage $1$, and $Q$ is
the full future space in stage $3$.
For a physical bank entry $d$, write $t=t_{d_j}$.

The stage expands its current tensor factor from a triple line to
all of $D$, while retaining the future line $Q$.  More precisely,
its incoming labels are
$(E\otimes\langle t\rangle)\otimes Q$ on $X$ and
$(B\otimes\langle t\rangle)\otimes Q$ on $Y$; its outgoing
labels are $(E\otimes D)\otimes Q$ on $X$ and
\[
 \bigl((B\otimes D)\mathbin\perp(P\otimes t^\perp)\bigr)\otimes Q
 = (P\otimes\langle t\rangle)^\perp_{E\otimes D}\otimes Q
\]
on $Y$.  Separating the next triple from $Q$ makes these exactly
the next stage's incoming labels.  At the first and last stages
they agree with the specified terminal labels.  We may insert these
stage boundaries as extra vertices without changing any edge operation.

The following gate labels realize those boundary conditions.  In the
table each label is further tensored with $Q$, and $t$ is the
indicator at the physical bank entry of an entry-specific gate.
The $X$ labels increase to $E\otimes D$, while the $Y$ labels
increase through $B\otimes D$.  A central wire must visit both
levels: its return from $E\otimes D$ in row $3$ to $B\otimes D$
in row $4$ is the only decrease that will consume part of the margin.

\begin{equation}\label{net:labels}
\begin{array}{c|l|l}
 \text{row}&\text{physical wires touched}&
                 \text{label before tensoring with }Q\\\hline
 0&Y,\ \text{side}& B\otimes\langle t\rangle\\
 1&Y,\ \text{center}& B\otimes D\\
 2&X,\ \text{side}&(B\otimes D)\mathbin\perp(P\otimes\langle t\rangle)\\
 3&X,\ \text{center}&E\otimes D\\
 4&Y,\ \text{center}&B\otimes D\\
 5&Y,\ \text{side}&(B\otimes D)\mathbin\perp(P\otimes t^\perp)\\
 6&X,\ \text{center}&E\otimes D\\
 7&X,\ \text{side}&E\otimes D.
\end{array}
\end{equation}
The row numbers here are chronological, including in stage $2$.
At that stage, reverse the eight scalar updates and their signs, with
logical source $Y$ and target $X$.  The physical touch sequence is
$Y$--side, $Y$--center, $X$--side, $X$--center, $Y$--center,
$Y$--side, $X$--center, $X$--side, exactly as in the table.
Relabeling the ordered side pairs accordingly leaves their two
physical triples neighboring.  Every displayed label is nondegenerate.

\begin{lemma}\label{net:residuals}
Every edge has nested nondegenerate endpoint labels.  Each nonzero
orthogonal residual, meaning the complement of the smaller label
inside the larger, has an orthonormal binary basis.  If $s$ is the
sum of the residual dimensions, then
\begin{align}
 s&=Wm-2N_0+2I_0(101)100
   =1873807244636671267308000000,\label{net:s}\\
 \Delta:=Wm-s&=6871402692000000>0.\label{net:delta}
\end{align}
\end{lemma}

\begin{proof}
Here is the full list of nonzero or potentially nonzero residuals.
All rows except the last are tensored with $Q$.  The vectors
$t_X,t_Y$ on a side wire are its two physical triple indicators.
The omitted edges have zero residual.
\[
\begin{array}{l|l}
 \text{wire and edge}&\text{orthogonal residual}\\\hline
 X:\ \mathrm{in}\to2&B\otimes t_X^\perp\\
 X:\ 2\to3&P\otimes t_X^\perp\\
 Y:\ 0\to1&B\otimes t_Y^\perp\\
 Y:\ 4\to5&P\otimes t_Y^\perp\\
 \text{center}:\ \mathrm{source}\to1&B\otimes D\\
 \text{center}:\ 1\to3,\ 3\to4,\ 4\to6&P\otimes D\quad\text{(each)}\\
 \text{side}:\ \mathrm{source}\to0&B\otimes\langle t_Y\rangle\\
 \text{side}:\ 0\to2&(B\otimes t_Y^\perp)\mathbin\perp
                              (P\otimes\langle t_X\rangle)\\
 \text{side}:\ 2\to5&P\otimes\langle t_X,t_Y\rangle^\perp\\
 \text{side}:\ 5\to7&P\otimes\langle t_Y\rangle\\
 \text{auxiliary}:\ \mathrm{last}\to\mathrm{sink}
                         &E\otimes D\otimes Q^\perp.
\end{array}
\]
In the last row $Q^\perp$ is taken in the full future tensor space.
The side edge $2\to5$ explains the use of neighboring triples.
Because $t_X\cdot t_Y=0$ and both vectors have norm one,
\[
 t_Y^\perp=\langle t_X\rangle\mathbin\perp
                    \langle t_X,t_Y\rangle^\perp.
\]
Its row-$2$ label $(B\otimes D)\mathbin\perp
(P\otimes\langle t_X\rangle)$ therefore lies in its row-$5$
label $(B\otimes D)\mathbin\perp(P\otimes t_Y^\perp)$,
with precisely the displayed residual.  The other rows follow by
$E=B\mathbin\perp P$ and
$D=\langle t\rangle\mathbin\perp t^\perp$.
These decompositions verify all inclusions and show that the only
decreases are the central edges $3\to4$, each of dimension $h$.

The signed changes telescope along physical wires.  Source
dimensions sum to $N_0$ and sink dimensions to $Wm-N_0$, so their
total is $Wm-2N_0$.  Each of the $I_0(101)$ central wires loses
$h=100$ dimensions once.  Replacing signed changes by absolute
changes adds twice these losses, proving the formula for $s$.
In particular,
\[
 \Delta=2N_0-2I_0(h+1)h
       =2v^2\bigl(v-3h(h+1)\bigr)>0,
\]
since $v>3h(h+1)$.  Substituting $h=100$ and $v=161700$ gives
the displayed integers.

Nondegeneracy alone would not guarantee an orthonormal basis in
characteristic two: a nondegenerate form may be alternating.  Here
every nonzero residual contains a vector of norm one.  Indeed,
$t^\perp$ and $\langle t_X,t_Y\rangle^\perp$ contain coordinate
units outside supports of size at most $3$ and $6$ respectively.
If nonzero, $B$ contains a coordinate unit outside a support of
size $3^{j-1}<h^{j-1}$, and the same argument applies to nonzero
$Q^\perp$.  Full spaces and the lines $P,Q,\langle t\rangle$
also contain norm-one vectors.  Taking tensor products and choosing
a nonzero summand in an orthogonal sum proves the claim for every
residual in the table.

For completeness, a nondegenerate binary symmetric space with a
norm-one vector has an orthonormal basis.  Split off unit lines
until the remaining space is alternating.  A nonzero nondegenerate
alternating space splits into orthogonal planes with bases $p,q$
satisfying $p\cdot p=q\cdot q=0$ and $p\cdot q=1$: choose $p\ne0$,
choose $q$ by nondegeneracy, and split off their plane.  Retain a
unit vector $w$ from a line already removed.  A plane perpendicular
to $w$ is absorbed by replacing $(w,p,q)$ with
\[
 w+p,\qquad w+q,\qquad w+p+q.
\]
These are pairwise orthogonal unit vectors and span the same
three-dimensional space.  Repetition absorbs all alternating
planes.  All choices can be made by finite binary linear algebra
with a fixed ordering of coordinates and pivots.
\end{proof}

\subsection{Transforming every physical role}

\begin{proposition}[Finite network interface]\label{net:finite-interface}
For every integer $f\ge1$, the explicit network applies
$C^{\otimes mf}$ to each of $W$ arbitrary arrays indexed by
$\Ftwo^{mf}$.  It uses $s$ directional steps, each on one physical
array and conjugate by address permutations to copies of
$C^{\otimes f}$ on fibers.  All other work consists of a fixed
number of pointwise linear maps and address permutations.  The
same assertion holds with any spectator address bits.
\end{proposition}

\begin{proof}
Use the gate and terminal labels just constructed in
\eqref{net:telescoping}, with each array address viewed as $f$ columns
of $m$ bits.  It remains to calculate the terminal frame ratios and
implement the edge changes.

Except for the route $X_d\to Y_d$, a source and its routed sink
have labels $0$ and $\mathcal D$, so their frame ratio is
$C^{\otimes mf}$.  For the exceptional route put $u=u_d$.
Since $u\cdot u=1$, $P_{U_d}x=u[u\cdot x]$, and orthogonal
weight additivity gives
\[
 q_{U_d^\perp}(x)-q_{U_d}(x)
 =\wt(x)-2\wt(u)[u\cdot x]
 =\wt(x)+2[u\cdot x]\pmod4.
\]
Here $\wt(u)=27$ is odd.  Hence the exceptional frame ratio is
$R_uC^{\otimes m}$ in each column.  Correct these translations
at the sinks and undo the signed bank exchange.  The result is
$C^{\otimes mf}$ on each physical role, including every
auxiliary array.

By Lemma~\ref{net:edge}, an edge with residual dimension $\rho$
uses $\rho$ products $C_z^{\otimes f}$ or
$(C_z^{-1})^{\otimes f}$.  Extend $z$ to a binary basis and apply
the corresponding invertible address change in each column.
Translation by $z$ becomes a flip of one selected bit, so the
forward product is $C^{\otimes f}$ on those $f$ selected bits,
for each setting of all other bits.  Inverse products add only
translations.  Summing the residual dimensions gives precisely
$s$ directional steps.  The network, all its bases, and the
number of remaining pointwise operations are fixed independently
of $f$.  Spectator bits merely index additional copies of these
identities.
\end{proof}

\subsection{Linear-time address changes}

The address permutations in Proposition~\ref{net:finite-interface}
must not introduce a factor $f$.  The following enumeration also
explains the uniform construction of the recursive batches.

\begin{lemma}\label{net:indexing}
Let $k=mf+r$ with $0\le r<m$, and fix an invertible binary map
on an $m$-bit column.  Applying it in all $f$ columns, translating
columns by prescribed fixed vectors, and collecting one chosen
output bit from every column into a contiguous $f$-bit address
field can be implemented in $O(2^k)$ time.  This includes all
index preparation and data movement; the constant may depend on
$m$ and the fixed maps, but not on $f$ or $k$.
\end{lemma}

\begin{proof}
An $m$-bit column is a digit of radix $2^m$.  Since $m$ and the
maps are fixed, a finite table gives, for each possible column
digit, its old bits, transformed bits, selected output bit, and
remaining output bits.  This table can be constructed once by
the specified binary arithmetic; its cost is an absolute constant.
The $r$ spectator bits form one additional digit of radix $2^r$,
omitted when $r=0$.

For a given $k$, compute the input and output bit-position strides
and the displaced table contributions for each column.  This
costs $O(f+1)$ operations with a constant depending on $m$.
An input address and the desired output address are each sums
of the contributions of their column digits and the spectator
digit; the output contributions put the selected bits in their
common field and all other bits outside it.  Traverse the digit
prefix tree, retaining these two partial sums.  Extending a prefix
by a digit adds its two tabulated contributions in constant time.
At a leaf, copy the value between the two resulting addresses,
using a separate output array to avoid overwriting unread values.
Every internal node has at least two children.  There are $2^k$
leaves, so fewer than $2^{k+1}$ nodes, including the root.
The preparation cost is also $O(2^k)$.  Reversing the source and
destination addresses implements the inverse permutation with
the same bound.

Take the selected $f$ bits as the least significant output field.
Each fiber is then contiguous, and consecutive groups of $W_*$
fibers, for any fixed $W_*$, are consecutive blocks of
$W_*2^f$ entries.  No per-entry scan of the columns is needed
to form or undo these groups.  Addresses have $k+O(1)$ bits,
with the fixed network sizes included in the constant.
\end{proof}

\subsection{Exact batching and the critical exponent}

The least power of two at least $W$ is
\begin{equation}\label{net:padding}
 W_*=2^{71},\qquad
 S=s+(W_*-W)m=W_*m-\Delta,\qquad
 \lambda=\frac{S}{W_*}=m-\frac{\Delta}{2^{71}},\qquad
 \theta=\log_m\lambda.
\end{equation}
In particular $1<\lambda<m$, so $0<\theta<1$.
The $W_*-W$ extra roles each use the $m$ ordinary coordinate
directions, preserving the absolute saving $\Delta$.

\begin{samepage}
\begin{theorem}\label{net:tensor-bound}
There is a deterministic exact algorithm applying $C^{\otimes k}$
to an arbitrary complex array of length $2^k$ in
\[
 O\bigl(2^k(k+1)^\theta\bigr)
\]
time, for every integer $k\ge0$.  The bound includes fixed-table
construction, all scalar arithmetic, array initialization, index
preparation, and data movement.  Only the rational constants and
$i$ are required.
\end{theorem}
\end{samepage}

\begin{proof}
We first transform $W_*$ arbitrary arrays simultaneously.  Set
$K=m(71+1)$.  For $k<K$, apply the $k$ factors of $C$ one at a
time along their fibers.  Because $K$ is fixed, the cost divided
by $W_*2^k$ is bounded by an absolute constant throughout this
base range, including $k=0$.

For $k\ge K$, write $k=mf+r$, where $f=\lfloor k/m\rfloor$
and $0\le r<m$.  Apply the $r$ remaining bit factors directly.
On $W$ arrays use Proposition~\ref{net:finite-interface} with
$f$ columns and $r$ spectator bits.  On each of the $W_*-W$
remaining arrays use one directional step for each of the $m$
coordinate positions, again simultaneously over all $f$ columns.
There are $S$ directional steps in total.

Each such step requires $2^{k-f}$ independent transforms
$C^{\otimes f}$.  Since $f\ge72$ and
$k-f\ge(m-1)f\ge71$, the number $2^{k-f}$ is divisible by
$W_*=2^{71}$.  After the permutation of
Lemma~\ref{net:indexing}, partition its fibers into exactly
$2^{k-f}/W_*$ batches and use the present $W_*$-array algorithm
recursively on each batch.  No partially filled recursive batch
occurs.  Its correctness follows by induction on $k$, since
$f<k$, and because the network interface applies to arbitrary
values on \emph{all} its roles.

Let $T(k)$ be the cost for $W_*$ arrays and
$t(k)=T(k)/(W_*2^k)$.  The fixed number of pointwise gates,
permutations, terminal corrections, and at most $m-1$ leftover
bit layers cost at most $A W_*2^k$ outside recursive calls,
for an absolute constant $A$.  Batch setup and returning results
are included in this bound: the batches partition a packed array,
and their number is at most its length.  Thus
\begin{align*}
 T(k)&\le A W_*2^k+
          S\frac{2^{k-f}}{W_*}T(f),\\
 t(k)&\le A+\lambda t(\lfloor k/m\rfloor)
                   \qquad(k\ge K).
\end{align*}
This exact recurrence permits the critical exponent
$\theta=\log_m\lambda$.  To see this directly, unroll it for
$d$ steps until the argument falls below $K$.  Successive
arguments are $\lfloor k/m^j\rfloor$, and
$d\le\lceil\log_m(k+1)\rceil$.  If $B$ bounds the normalized
base costs, then
\[
 t(k)\le A\sum_{j=0}^{d-1}\lambda^j+B\lambda^d
       \le\left(\frac{A}{\lambda-1}+B\right)\lambda^d
       =O((k+1)^\theta).
\]
Here $\lambda^d\le\lambda(k+1)^\theta$.  The same estimate
holds in the base range by increasing the constant.

Fixed tables may be generated once and reused; even generating
them on the first invocation contributes an absolute constant.
The recursive calls can run serially and reuse their working
arrays.  Their address fields, counters, and strides use
$O(k+1)$ bits: at each active level only a fixed number of
arrays of the current length is needed, and the lengths shrink
from $2^k$ to $2^{\lfloor k/m\rfloor}$ and so on.
The recursion stack has $O(\log(k+2))$ levels.  Thus the claimed
workspace is bounded by a constant times
$\sum_{j=0}^{d}W_*2^{\lfloor k/m^j\rfloor}+O(k+1)
=O(W_*2^k)$; consecutive nonterminal array lengths decrease
by at least a factor of two.  The claimed index costs therefore
use exactly the logarithmic-word model of the paper.

Finally, initialize $W_*-1$ unused arrays to zero, put the
requested array in the remaining role, run the simultaneous
procedure, and retain that role's output.  Initialization costs
$O(W_*2^k)=O(2^k)$, and $W_*$ is an absolute constant.  This
proves the stated single-array bound.
\end{proof}

\begin{remark}\label{net:role-remark}
Restoring arbitrary auxiliary values in Lemma~\ref{net:scalar}
is essential.  After the frame changes, those physical roles
are required to undergo the same tensor transform as the data
banks.  A network valid only with zero auxiliary values would
not justify using arbitrary transform fibers as all $W_*$ roles
in the recursive batches.  Equation~\eqref{net:telescoping}
establishes the required full linear-map identity.
\end{remark}

